\chapter{Christ's Witnesses at Rome}
"A great multitude."--\textit{Tacitus}.


With jubilation Quintus sees again the shores of Italy rise over
the Adriatic, and finds himself once more in his beloved Rome.  The
center of magnificence and power it seems.  Alter clamorous public
greetings in the Forum, there comes another welcome which happens
only in a returning soldier's life.  In the palace of Marcus the
kindred of Quintus are gathered, and Lucretia also is in the
circle, to hear his great adventure.

"How wonderful it seems," the knight begins: "so many times have
your faces come to me in my dreams, but now I am fully awake and
see them once again.  Hail to you all!  When I was sailing away
from Brundisium, the augur foretold for me an unusual experience.
In the Jewish life beyond the Sea I have learned much, if that were
the fulfillment.  But, most of all, I have come back with a new
religious faith.  In Judaea, as you must have heard, a certain
Galilaean has called himself the Son of the one true God.  He has
spoken of a future life for men; and he has now risen from the
grave, after his torture on a cross, to prove his doctrine true.  I
now believe in him, as the interpreter of the future life.
Forevermore he is my High Priest, and not the great pontifex in the
temple of your Jupiter."

Brave words they are.  There in the great hall of Marcus, with the
sunlight shining on the gorgeous palaces of the Caesars, the Temple
of Apollo, and all else which crowns the Palatine, the noble
Quintus confesses his new belief.  Come what will the consequences!

Then, while they hear in amazement, he further says; "Most inviting
is this new creed.  Our wise Roman scholars, as well as those in
Greece, have only been guessers about the future life.  But the
Christus speaks as one who has come from the heavens.  Those who
keep his commandments are to dwell with him forevermore in eternal
joy.  Everywhere through Judaea men are becoming his followers, and
the wide world is to believe on him.  Perhaps you also, my
cherished ones, will come to accept his teaching of the future
life."

So Quintus speaks, with his vibrant voice and with a strange light
on his face.  Wonderingly they hear the tidings that he brings--the
recital of the greatest happening that can ever befall a man.  Not
deriding their valiant soldier, and not withholding their wealth of
love from one who has come safely back to them, they watch the
changes in his life.

"I do not care," he says, "to loiter in the baths of Agrippa and to
hear from the idlers there the gossip of the hour.  The
gladiatorial struggles in the Circus Maximus and the comedies in
the theaters have lost for me their relish.  For the civic rewards
which Tiberius gives his favored ones I have no wish.  Senatorships
and proconsulships are like the dust in the apothecaries' scales.
I have seen the risen Lord!"

Influential is such a life on the home group of Quintus.  With his
pride of birth and his great properties, Marcus becomes a believer.
A conversion it is which is the surprise of Rome.  The rare
Lucretia, as well, receives the truth.  At times, before she has
called herself a disciple, Quintus escorts her to the worship of
the Roman Christians.  Their captivating speech, their holy love
for one another, their rapturous faces move her deepest heart.
Till, one day, when Quintus has been telling her of the womanhood
in Judaea which the Christ has ennobled, she replies:

"I believe it all, O Quintus.  Of late into my heart an untold
peace has come.  All things are changed for me.  The sunlight is on
the hills!"  It is her open confession.  Lucretia is thenceforth
enrolled among the Roman saints of whom the world was not worthy,
and who looked for the life to come.


In the fellowship of the Roman church--already founded and rapidly
enlarging--Quintus finds his pleasure.  A few are Jews from the
ghetto beyond the Tiber, till the persecution of Claudius drives
them forth.  More are of the varied nationalities met in that
commercial and luxurious center.  Most are of plebeian blood.
There are smiths and mechanics; there are stone cutters, workers in
mosaics, and decorators.  There are slaves from the very palace of
Tiberius.  There is Amon from Egypt, who sells his jewelry down in
the Nova Via.  There is Polemon, the Grecian shopkeeper, in the
Clivus Victoriae.  There is Onesimus, the servant of Philemon, from
Colossae.  There are Amplias and Epaenetus and Stachys, the
particular friends of the Gentile apostle.  There is, as well,
Pomponia Graecina, that woman of noble blood, who accepts the
Christ.  An ever-increasing company it is.

In their assemblies, on the first day of the week, Quintus has his
influential place.  He listens to the reading of the older
Scriptures; he celebrates with the gathered company the eucharistic
suppers and agapae; he keeps with them the Easter celebration, in
memory of Him who shall give them eternal life.  In emblem of their
faith the sign of the fish is on their evening lamps.  Theirs is a
sterling citizenship.  The wanton metropolis of the Caesars is
blessed immeasurably by the company of these who follow the risen
Lord.

It is after the midcentury that the great Paulus, having met with
shipwreck on Melita, draws near to Rome.  Quintus leads the company
that goes out southward forty miles, to welcome the Christian
traveler.   At Appii Forum, that common town with its bargemen and
its tavern keepers, they give the kiss of welcome to a little bent
and gray-haired Jew, who shall go down into history as Christ's
most illustrious apostle.  The faithful Luke is his companion.
Along the famous highway of the Via Appia, where emperors and
warriors, scholars and Oriental tradesmen have walked, Quintus
escorts their guest.  Past the tombs of the Roman great, by
uncounted statues, past suburban villas they go, until, through the
Porta Appia, the holy prisoner, chained to a Roman guard, finds
himself in the city of the Caesars.

One rare privilege the Roman knight then envoys.  In his hired
house, near the Pretorian camp, Paul speaks without interruption
his words of grace.  The doctrines he had before written to the
Roman church he now explains; the wish he had made to see them face
to face now expresses itself in words of love.  The flood tides of
his eloquence move resistlessly on, as he interprets the new faith
and speaks of Him who is to give them eternal life.  Quintus is
enriched by his frequent association with the peerless soul.  Nor
did he have a prouder thing to say, in the days to come, than to
declare, "I heard great Paulus tell of the life immortal."

But how fares our knight when persecution comes?  Through the years
he has been bravely declaring the Christian doctrine of the eternal
life to priests in the temples, to Roman nobles, to all most
hostile.  But his wealth and social standing, as well as the
emperor's favor, now insure his safety.  His father Marcus has long
since passed on, in hope of the heavenly life.  Having wedded the
graceful Lucretia, when an apostle was in Rome to speak their
nuptials, he has her efficient counsel in the testing times.

"Look! look! Lucretia," he cries, one evening; "through the lower
city the flames are running like unbridled horses.  There is danger
that all Rome may go to ashes."

For nine long days they watch the sweep of the lurid flames.  The
light shines out like a signal torch, to mark an emperor's folly.
Then the undeserved charge that they have lit the flames brings on
the martyrdom of the Roman Christians.  Sometimes Quintus and
Lucretia are able to soften the trials of the sufferers, by
permission of the capricious Nero.  To old Chilo, the Grecian,
before he meets his doom, they unfold the promise of eternal reward
in the Father's house.  The hope of immortality they carry to those
who go to the lions, at the emperor's whimsical command.  And the
glorious company of martyrs passes singing to the skies, because of
their consoling words.

Down into the dungeon of the Mamertine they are permitted once to
go, to visit Paulus.  But he needs not their consolation.  Rather
he is the comforter.  With the poise of a conqueror he bids them
not to mourn for him: he is going to the Lord in the unending life.
Over their bowed heads he stretches his aged hands, in apostolic
benediction.  Soon ends his imprisonment.  At \textit{Tre Fontane}, in a
few days more, his weary body rests; but his immortal spirit mounts
beyond the stars.

At last the Christian knight comes to the crossing.  The prediction
of the augur at Brundisium has been strikingly fulfilled.  Matured
in all the graces, he is like the ripened Chian clusters that await
the vintager in the autumn days.  The friends of Quintus have gone
before; as the old century wanes, the old man is to follow them.

"My time has come to go," he says one day; "the portals of eternal
life and joy I see swinging open wide.  I shall pass through the
gates, because my ascended Lord has gone in before me to prepare my
dwelling place.  With him as my Teacher I believe in the life
immortal."


In the Roman catacombs, those most remarkable testimonies to the
eternal life, his resting place may be found.  The sign of the fish
is on his stone.  Its time-eaten inscription is still legible,
among the many which tell of the early Christian expectation and of
all future Christian hope:


"HERE RESTS THE DUST OF QUINTUS, OF NOBLE BLOOD; IN THE FAITH OF
THE ASCENDED LORD HE HAS ENTERED UPON THE ETERNAL LIFE."
